\documentclass{article}
\usepackage{graphicx} 
\usepackage{amsmath}
\usepackage{enumitem}
\usepackage{amsthm}
\newtheorem{theorem}{Theorem}
\usepackage{multirow}
\usepackage[margin=1in]{geometry}
% \renewcommand{\thetheorem}{\arabic{theorem}.}

\title{Lab 1 - Technical Communication}
\author{UTKARSH KUMAR 23114101}
\date{January 2026}

\begin{document}
\setlength{\parindent}{0pt}
\maketitle

\begin{enumerate}

\item

Let $h \in \mathbf{H}$. Define a indicator random variable $X_{i,j}$ such that
\begin{center}
$X_{i,j}$ = \begin{cases}
    \text{1,} & \text{if $h(k_{i})=h(k_{j})$} \\
    \text{0,} & \text{o.w.}
\end{cases}
\end{center}
Let the length of chain at index $h(k_{i})$ be a random variable $X_{i}$. Then we have
\begin{center}
    $X_{i}$ = $\sum\limits_{j=0}^{n-1} X_{i,j}$
\end{center}
The expected length of chain at index $h(k_{i})$ is then given by
\begin{align*}
    \underset{h\in\mathbf{H}}{E}\bigl(X_{i}\bigr) &= \underset{h\in\mathbf{H}}{E}\bigl(\sum\limits_{j=0}^{n-1} X_{i,j}\bigr)\\
    &= 1 + \ \underset{h\in\mathbf{H}}{E}\bigl(\sum\limits_{j=0,j\not=i}^{n-1} X_{i,j}\bigr)\\
    &= 1 + \sum\limits_{j=0,j\not=i}^{n-1} 1 \cdot \underset{h\in\mathbf{H}}{\Pr}(h(k_{i})=h(k_{j}))+0 \cdot\underset{h\in\mathbf{H}}{\Pr}(h(k_{i})\not=h(k_{j}))\\
    &= 1+ \sum\limits_{j=0,j\not=i}^{n-1} \frac{1}{m}\\
    &= 1+\frac{n-1}{m}
\end{align*}
Since $m=\Theta(n) \implies \alpha = O(1) \implies$ expected length of chain at index $h(k_{i})$ is constant.

\item
The construction of the \textbf{dot product} family of hash function is defined as follows.

\begin{enumerate}[label=(\alph*), leftmargin=2.5em]
    \item Let $|U| = m^r$ for some prime $m$and positive integer $r \ge2$.
    \item For a key $k \in U$, view $k$ in base $m$ representation. So, we have 
        \begin{center}
            $k = (k_{0},k_{1},...,k_{r-1})$
        \end{center}
    where $k_{i} \in \{0,1,...,m-1\}$.
    \item Let $a = (a_{0},a_{1},...,a_{r-1})$ where $a_{i}$'s are picked uniformly at random from $\{0,1,...,m-1\}$.
    \item Define $h_{a}(k):U \rightarrow \{0,1,...,m-1\}$ such that
        \begin{center}
            $h_a(k) = \sum\limits_{i=0}^{r-1} a_{i}k_{i} \bmod m.$
        \end{center}
    \item $\mathbf{H} = \{h_{a}\}$
\end{enumerate}

\item
\begin{theorem}
    The probability of a collision in the dot product family of hash functions is $\frac{1}{m}$.
\end{theorem}

\begin{proof}
    Let $x = \{x_{0},x_{1},...,x_{r-1}\}$ and $y = \{y_{0},y_{1},...,y_{r-1}\}$ be two keys such that $x \not = y$. Then they differ in at least one digit position. Without loss of generality assume that $x_0 \not = y_0$. Then we have

    \begin{align*}
        \underset{h\in\mathbf{H}}{\Pr}(h(x)=h(y)) &= \underset{h\in\mathbf{H}}{\Pr}\bigl(\sum\limits_{i=0}^{r-1} a_{i}x_{i}\bmod  m=\sum\limits_{i=0}^{r-1} a_{i}y_{i} \bmod m\bigr)\\
        &= \underset{h\in\mathbf{H}}{\Pr}\bigl(\sum\limits_{i=0}^{r-1} a_{i}x_{i}\ \equiv \sum\limits_{i=0}^{r-1} a_{i}y_{i} \bmod m\bigr)\\
        &= \underset{h\in\mathbf{H}}{\Pr}\bigl(a_0(x_0-y_0) + \sum\limits_{i=1}^{r-1} a_{i}(x_{i}-y_i)\ \equiv 0 \bmod m\bigr)\\
        &= \underset{h\in\mathbf{H}}{\Pr}\bigl(a_0 \equiv (x_0-y_0)^{-1} \cdot \sum\limits_{i=1}^{r-1} a_{i}(x_{i}-y_i)\bmod{m}\bigr)\\
        &= \frac{1}{m}
    \end{align*}

    Since $x_0 \not = y_0$ and $m$ is prime, inverse of $x_0-y_0$ exists.
\end{proof}

\item

In Table 1. we show how to draw a table with \LaTeX.

\begin{table}[h]
\centering
\renewcommand{\arraystretch}{1.2}
\begin{tabular}{c|ccccc}
\hline
\multicolumn{1}{c|}{Row $i$} & \multicolumn{5}{c}{Columns}\\
\cline{2-6}
\multicolumn{1}{c|}{} & $X^i_0$& $X^i_1$& $X^i_2$& $X^i_3$& $X^i_4$\\
\hline
2&2&1&1&2&2\\
3&3&3&4&4&3\\
4&7&8&7&7&6\\
5&8&8&13&14&15\\
6&45&45&32&31&29\\
\hline
\end{tabular}
\caption{Demo of the table.}
\end{table}
\end{enumerate}
\end{document}
